\documentclass[11pt]{article}
\usepackage[margin=1in]{geometry}
\usepackage{hyperref}
\usepackage{enumitem}
\usepackage{booktabs}
\usepackage{xcolor}
\usepackage{amsmath}
\title{ProbOS --- Month 3, Week 11\\Pharmaceutical Stability Model}
\author{Nisong Monyimba}
\date{\today}

\begin{document}
\maketitle

\section*{Week 11 goal}
Build \texttt{PharmaStabilityModel}, reusing
\texttt{BatteryModel2Cell}'s Arrhenius reaction-kinetics architecture
with drug potency as the depleting quantity, per the near-term
priority established in \texttt{docs/monthly\_plans/month3/main.tex}
and the strategic reasoning in \texttt{docs/vision/main.tex}.

\textbf{Critical constraint, stated up front:} this model's
credibility rests entirely on being validated against a REAL,
genuine published stability-study dataset -- the same standard
Kim et al. (2007) set for \texttt{BatteryModel2Cell} in Month 1
Week 2. Synthetic or fabricated validation data would violate this
project's established honesty discipline
(\texttt{docs/standards/quality\_standards.md}) and must not be used.
If a suitable public dataset cannot be found by the end of Day 1,
this week's scope will be explicitly narrowed (e.g. to model
architecture and unit correctness only, with validation flagged as a
genuinely open item) rather than proceeding with fabricated data.

\section*{Day-by-day plan}

\subsection*{Day 1 -- Literature research (before any model code)}
\begin{itemize}[leftmargin=*]
  \item Search for a real, published pharmaceutical stability study
    reporting accelerated and/or long-term degradation data (ICH
    Q1A-style: potency vs. time at multiple temperatures), analogous
    to how Kim (2007)'s ARC data was used for the battery model
  \item Confirm the underlying degradation mechanism is genuinely
    Arrhenius (most small-molecule drug degradation is, but this must
    be confirmed for the SPECIFIC dataset chosen, not assumed)
  \item Document the citation precisely (matching the existing
    standard: exact source, exact reported values, no invented
    numbers) before writing any model code
\end{itemize}

\subsection*{Day 2 -- PharmaStabilityModel implementation}
\begin{itemize}[leftmargin=*]
  \item Design the state vector and parameter vector based on the
    ACTUAL structure of the dataset found on Day 1, not assumed in
    advance -- unlike the battery model's 2-cell, 3-reaction
    structure, a real stability study may report a single degradation
    pathway, requiring a simpler state vector
  \item Implement \texttt{forward\_batch()} following the exact same
    vectorised-NumPy, CPU/GPU dispatch-ready pattern established in
    \texttt{BatteryModel2Cell} (Week 10's \texttt{xp} dispatch and
    \texttt{cp.fuse()} pattern should be reused directly, not
    reinvented)
\end{itemize}

\subsection*{Day 3 -- Validation against the real dataset}
\begin{itemize}[leftmargin=*]
  \item Deterministic validation: run the model at the dataset's own
    reported nominal parameters, confirm it reproduces the reported
    potency-vs-time curve within a stated, justified tolerance --
    matching the standard already set for Kim (2007)
  \item If the model does NOT match the real data, this must be
    reported honestly and investigated, not hidden or the tolerance
    quietly loosened until it passes
\end{itemize}

\subsection*{Day 4 -- Monte Carlo + Sobol sensitivity}
\begin{itemize}[leftmargin=*]
  \item Build parameter priors reflecting genuine
    manufacturing/formulation uncertainty (informed by the dataset's
    own reported variability where available, not invented)
  \item Run \texttt{SobolSensitivity} to identify which parameters
    drive shelf-life uncertainty -- the direct pharma analogue of
    \texttt{Ea\_SEI}'s dominance in the battery model
\end{itemize}

\subsection*{Day 5 -- Test coverage}
\begin{itemize}[leftmargin=*]
  \item \texttt{python/tests/test\_pharma\_stability.py}: shape
    contracts, physical invariants (potency cannot exceed 100\% or go
    negative), and the Day 3 validation-against-real-data test as the
    centerpiece, matching \texttt{test\_battery\_model.py}'s standard
\end{itemize}

\subsection*{Day 6 -- pre\_week\_audit.sh + CI}
\begin{itemize}[leftmargin=*]
  \item Run \texttt{pre\_week\_audit.sh} before considering the week
    complete
  \item Add the new model to CI's example-running steps
\end{itemize}

\subsection*{Day 7 -- Documentation + retrospective}
\begin{itemize}[leftmargin=*]
  \item Update \texttt{docs/vision/main.tex}: move pharmaceutical
    stability from ``Structurally ready'' to ``Validated'' IF Day 3's
    validation genuinely succeeds against real data -- if it does
    not, leave the honest status as-is and document why
  \item Retrospective appended to this document once complete
\end{itemize}

\section*{Exit criteria}
\texttt{PharmaStabilityModel} exists, is validated against a real,
cited, published stability dataset (not fabricated data) at an
honestly-justified tolerance, and its dominant uncertainty driver is
identified via Sobol sensitivity -- matching the exact validation
discipline already established for \texttt{BatteryModel2Cell}. If a
suitable real dataset cannot be found, this is reported honestly as a
genuinely open item, not worked around with synthetic data presented
as if it were real.


\section*{What was actually accomplished (retrospective)}

\begin{itemize}[leftmargin=*]
  \item \texttt{PharmaStabilityModel} built and validated against a
    REAL, peer-reviewed, ICH-referenced dataset: Gonzalez-Gonzalez et
    al. (\emph{Molecules} 2023, 28(23), 7925), chlorhexidine stability
    data. Per the plan's own Day 1 gate, an initial candidate
    (aspirin hydrolysis kinetics from a 2024 teaching lab module) was
    explicitly rejected on closer scrutiny: real, peer-reviewed, and
    numerically correct, but pedagogical rather than an actual
    ICH Q1A-style product stability study.
  \item A genuine architectural finding emerged: the real,
    best-validated pharma dataset uses Avrami kinetics (a closed-form
    $\text{potency} = \exp(-k t^n)$ relationship), not the first-order
    depletion law used by \texttt{BatteryModel2Cell}. Code reuse
    happened at the architecture level (\texttt{Model} ABC,
    \texttt{xp} CPU/GPU dispatch, Sobol/provenance integration), not
    at the equation level.
  \item A genuine fitting failure was found and fixed: an
    unconstrained fit of both the Avrami pre-exponential factor (A)
    and exponent (n) against the paper's 3 real reported 365-day
    degradation percentages converged to a mathematically passing but
    physically absurd result ($A = 1.5 \times 10^{255}$,
    $n = -94.6$) -- a textbook degenerate/underdetermined fit. Fixed
    by fixing $n=1$ as a labeled simplifying assumption and fitting
    only $A$, which achieved an IDENTICAL residual error, confirming
    the extra degree of freedom was pure unidentifiable noise.
  \item The probabilistic core was reconnected to this new domain:
    using the paper's own reported Ea standard error ($\pm$2.61
    kcal/mol, genuine literature-backed uncertainty), Sobol
    sensitivity identified Ea as overwhelmingly dominant
    ($S_1 = 0.993$), and Monte Carlo propagation revealed a real tail
    risk: the worst-case 5\% of outcomes show complete potency loss
    within one year, versus 82.6\% potency at the point estimate.
\end{itemize}

\section*{What went right}
\begin{enumerate}[leftmargin=*]
  \item The Day 1 research gate worked as intended: scrutinizing the
    first candidate dataset caught a real quality gap before any
    model code was built around it.
  \item A direct question about whether the probabilistic core still
    mattered led to genuinely reconnecting the deterministic
    groundwork to real Monte Carlo + Sobol analysis, rather than
    treating the curve-fit as the finish line.
\end{enumerate}

\section*{What went wrong, and was fixed}
\begin{enumerate}[leftmargin=*]
  \item An unconstrained 2-parameter fit produced a passing residual
    error alongside physically nonsensical parameter values. Caught
    by inspecting the actual fitted values, not just the error
    metric, and fixed by removing the underlying degeneracy.
\end{enumerate}

\section*{Numbers at Week 11 close}
\begin{itemize}[leftmargin=*]
  \item Python tests: 398/398 passing (12 new)
  \item mypy strict: 0 errors
  \item ruff: 0 warnings
  \item Coverage: 89.18\%
  \item \texttt{pre\_week\_audit.sh}: 16/16 checks passed
  \item Validation accuracy: max 1.18 percentage points against 3
    real reported data points
  \item Sobol dominant parameter: Ea ($S_1 = 0.993$)
\end{itemize}

\section*{Carried into future weeks}
The Avrami exponent $n$ remains a fixed simplifying assumption, not
an independently recovered value. Automotive/EV safety positioning
(also a Month 3 near-term priority) remains untouched -- it requires
no new model, only outreach positioning around the already-validated
battery model.
\end{document}
